\section{Approach}\label{sec:method}

\subsection{Overview}

Our goal is to generate character animations that are both visually realistic and logically coherent with multimodal inputs. To achieve this, we introduce a framework designed to simulate both System 1 (reactive) and System 2 (deliberative) cognitive processes. Our model is built upon a Diffusion Transformer (DiT) backbone~\cite{seawead2025seaweed,gao2025seedance,esser2024sd3,dit}, which is first pre-trained~\cite{gao2025seedance} on general video generation tasks to acquire foundational video generation capabilities. We then transform this base model into a logical and expressive avatar through two critical designs, which are detailed in the following sections and illustrated in Figure~\ref{fig:framework}.

\par
Agentic Reasoning for Deliberative Control (Sec. 3.2): We first employ MLLM-based agents to reason about the input context and generate high-level semantic guidance. This step provides the deliberative control necessary to simulate ``System 2''. 

Multimodal Fusion for Reactive Rendering (Sec. 3.3): Next, our specialized MMDiT~\cite{esser2024sd3} architecture fuses this semantic guidance with reactive signals like audio  to simulate ``System 1''. To resolve modal conflicts, this architecture incorporates a pseudo-last-frame identity strategy, which prevents the static reference image from interfering with dynamic motion during training. 
\par
Beyond these core designs, our framework aligns with common practices~\cite{lin2025omnihuman,gan2025omniavatar,kong2025multitalk} for simplicity. To support long-form video synthesis, our framework can operate autoregressively, continuing generation by using the final frames of a previously generated clip as the initial frames for a new segment~\cite{stypulkowski2024diffused}. Our framework operates in the compact latent space of a pre-trained 3D VAE~\cite{yu20233DVAE} and is trained with a flow matching~\cite{liu2022reflow} objective. We omit further discussion of these standard components to focus on our main contributions.

\begin{figure}[t]
  \centering
  \includegraphics[width=\linewidth]{figs/framework.jpg}
  \caption{\textbf{The Dual-System Simulation Framework.} Our framework models avatar behavior by integrating a deliberative \textbf{System 2} for planning with a reactive \textbf{System 1} for rendering. \textbf{Left:} The overall pipeline. System 2 uses an MLLM Agent to reason over all inputs (audio, image, text) and generate a high-level "schedule". This schedule guides System 1's MMDiT network, which synthesizes the final video by fusing information within its dedicated text, audio, and video branches. \textbf{Right:} Key module details. (a) The reasoning pipeline consists of an MLLM Analyser and Planner that work together to create the schedule. (b,c) Our proposed \textbf{MM-Branch Warm-up} and \textbf{Pseudo Last Frame} methods mitigate multimodal conflicts during training.}
  \label{fig:framework}
\end{figure}


\subsection{Agentic Reasoning for Deliberative Control}



\myparagraph{Module Inputs and Outputs.} To model the deliberative nature of ``System 2'', our Agentic Reasoning module reasons over the input conditions to generate high-level, logically coherent guidance. The inputs are the character's reference image and the corresponding audio clip, supplemented by an optional text prompt describing the desired character behavior. The module is designed to process these inputs to produce semantic conditions in two forms: Reasoning Text, the explicit chain-of-thought text from the agents used directly as a condition, and Reasoning Latents, intermediate MLLM features extracted to serve as an additional conditioning signal and integrated via a dedicated attention mechanism. Our primary approach utilizes the former, while we have also investigated the integration of the latter.


\myparagraph{Multi-Step Reasoning Pipeline.} As shown in the top-right portion of Figure~\ref{fig:framework}, this deliberative guidance is generated using two MLLMs that play distinct roles in a collaborative process. The first MLLM, the Analyzer, receives the reference image, a corresponding text caption (produced by an auxiliary model to ensure accurate image interpretation), the audio clip, and an optional user-provided text prompt. Guided by a chain-of-thought prompt, the Analyzer MLLM performs an iterative reasoning process on the inputs. It systematically infers the character's persona, language style, speech content, emotion, intent, and environmental context, consolidating these insights into a single structured representation, typically a JSON object. This context-rich information is then passed to the second MLLM, the Planner. The Planner receives the Analyzer's output, along with the original character image for visual context. Guided by a new instructional prompt, its task is to devise an action plan structured as a sequence of shots, where each shot defines the character's expressions and actions for a duration corresponding to a single generation pass of our diffusion model. This collaborative reasoning yields a comprehensive motion schedule that preserves a coherent character persona through consistent actions across the entire video.

\myparagraph{Reflective Re-planning.} To maintain logical coherence in long-form video generation, our agentic framework incorporates an optional ``reflection'' process. During autoregressive synthesis, the generation schedule is dynamically updated by re-evaluating the most recently generated output. This process mitigates a common challenge in diffusion-based synthesis, where subtle execution deviations can accumulate and degrade logical coherence over time, particularly in longer videos.
In practice, the Planner takes the last generated frames and the original reference image as new inputs to re-assess its plan. This reflective loop corrects for semantic drift and helps maintain the video's logical consistency.


 
 \myparagraph{Investigation of Latent Feature Conditioning.} We also investigated the aforementioned approach of utilizing latent features from a MLLM as a semantic conditioning signal. The approach directly utilized the audio tokens within the transformer of the Analyzer agent. The rationale was that cross-modal attention in the transformer layers would enrich these audio token representations with high-level semantics (e.g., inferred emotion and intent) while preserving their original temporal structure. Based on this consideration, we selected these ``reasoning-infused'' audio latents from the final transformer layer and concatenated them with the raw audio features. This combined signal then replaced the original audio input for the DiT network.
\par

The above designs enable our agent to formulate a global, coherent plan for the entire scene. Unlike purely reactive methods that merely emulate ``System 1'', our approach further integrates deliberative reasoning from ``System 2'' to provide thoughtful, top-down guidance.

\subsection{Reactive Rendering via Multimodal Diffusion}


In this subsection, we describe how our diffusion model synthesizes the final video. It synergistically combines the high-level reasoning from the agents (primarily represented as text) with the low-level, reactive signals of audio inputs (primarily represented as audio features).


\myparagraph{Rethinking Reference Conditioning.} Before detailing our driving condition modeling, we must first analyze a critical input in video avatar models: the conditioning image, which serves two distinct purposes. The first is providing initial frames for autoregressive continuity, a standard practice we adopt by concatenating ground-truth (GT) frames. The second, more problematic purpose, is using a reference image for identity preservation. While early works used dedicated networks~\cite{aa,tian2024emo,zhu2023tryondiffusion} and recent methods reuse model parameters~\cite{lin2025omnihuman,kong2025multitalk}, both approaches inject a reference image sampled from the training video. As illustrated in Figure~\ref{fig:plf}, we argue that this creates a critical artifact: the model learns a spurious correlation that the reference image should appear literally within the generated sequence, severely restricting motion dynamics and conflicting with the audio and text driving signals. While this artifact can be partially mitigated by probabilistically sampling reference images from outside the training video clip, this approach may introduce new problems, causing the model to learn that the generated output should exhibit significant variation from the reference image. 

\begin{wrapfigure}{r}{0.5\textwidth}
    \centering
    \vspace{-12pt}
    \includegraphics[width=\textwidth]{figs/plf.jpg}
    \vspace{-18pt}
      \caption{\textbf{Rationale for the Pseudo Last Frame.} \textbf{Left:} Reference-conditioning has trended toward simplification. \textbf{Right:} The dilemma of reference image sampling. Sampling from \textcolor{green!50!black}{\textbf{within}} the target video segment ensures high relevance but restricts motion diversity. Conversely, sampling from \textcolor{magenta}{\textbf{outside}} the segment, a scenario that becomes more frequent as datasets grow larger and more dynamic, leads to a drop in content relevance, causing inconsistencies that undermine the reference's purpose.}
  \label{fig:plf}
\end{wrapfigure}



\par
The root cause of this issue is that the reference image is an artificial construct, not a condition native to the video data itself. It is for this reason that our solution is to discard it entirely during training and, in its place, introduce a novel guidance mechanism. As shown in the bottom right of Figure~\ref{fig:framework}, during training, we probabilistically condition the model on both the GT first and last frames of the video clip, as these are both native signals. During inference, we repurpose this mechanism by placing the user's reference image in the last frame's position, creating a pseudo last frame. Crucially, we shift its positional encoding (e.g., RoPE~\cite{rope}) to maintain a fixed temporal distance from the generated content. This pseudo-frame, which is dropped after rendering, functions as a ``carrot on a stick'': it guides the model toward the reference identity without ever forcing it to replicate the image. As our experiments show, this approach eliminates training artifacts and mitigates autoregressive error, achieving a superior trade-off between motion dynamics and stability.


\myparagraph{Symmetric Fusion and Warm-Up.} With all training conditions now native and compatible, we address the challenge of joint modeling. We adopt an MMDiT backbone but depart from prior work in our approach to audio conditioning. Instead of injecting audio features via additional cross-attention layers, we introduce a dedicated audio branch, architecturally symmetric to the video and text branches. All three modalities are then fused at each layer through a \textbf{shared} multi-head self-attention mechanism. This symmetric design offers two key advantages. First, it allows audio features to be iteratively refined alongside video and text, ensuring deep semantic alignment. Second, it enables true joint modeling, as tokens from all three modalities mutually attend to one another, facilitating a more effective mapping to a shared semantic space. While this new branch adds parameters, the computational overhead is negligible due to the low ratio of audio to video tokens.
\par
This symmetric architecture, while beneficial, introduces a training challenge. Naively training the entire model jointly leads to modality conflict, where the model learns to over-rely on the temporally dense audio signal for all predictions, thereby ignoring or washing out the more abstract guidance from the text branch. Freezing the pre-trained branches is also suboptimal, as it causes the audio branch to overfit and erroneously learn non-audio-related attributes like lighting and camera motion. To resolve this, we propose a two-stage warm-up strategy. In Stage 1, we train the full three-branch model jointly, forcing the model to learn an optimal division of labor: the text and video branches handle high-level semantics, compelling the audio branch to specialize in its core competencies (e.g., lip sync, speech mannerisms). For Stage 2, we construct the final model. The text and video branches are initialized with their original pre-trained weights, while the audio branch is initialized using the warmed-up weights obtained from Stage 1. This model is then fine-tuned. This strategy ensures each branch begins with its own strong, specialized prior, mitigating modality conflict and allowing each input to retain its distinct conditioning power. 
\par
Ultimately, the proposed architecture executes the deliberative plan. By redesigning the reference conditioning and equipping the model with an audio conditioning branch, our rendering process faithfully translates the high-level guidance from System 2 while maintaining the reactive fidelity of System 1.